% =====================================================================================================
% SGE001 -- fourth-order ("sharp") certificate of improvement.  Block ready to paste into manuscript/main.tex.
% Author of the block: theory agent, 3 October 2026.  Status of the mathematics: proved here (Taylor, assembled).
%
% WHERE TO PASTE
%   Right after \end{remark} of Remark~\ref{rem:orders} (end of the "Smooth regime" subsection) and before
%   \subsection{Hierarchy}.  Paragraph [E1-add] goes at the end of the paragraph "Bounds and certificates" of E1;
%   paragraph [E4-add] after the sentence that reports \EfourCertBoxFracPointFour in E4.  Everything uses the
%   macros of main.tex (\R, \Yh, \Hf, \status, \LN, \SU) and NEW numeric macros listed below, to be generated by
%   make_numbers.py from theory/sharp_certificate_results.json (no number below is typed by hand; the numbers
%   typed in the text are the orders 1..4, the factorials and the grid design).
%
% NEW MACROS  (S = res["E1"][<cell>]["summary"], P = res["prediction_LN_scaled"], F = res["E4_no_capacity"],
%              T = res["CES_enclosure_tightness"]; sigma macros with f"{x:.2g}", ratios f"{x:.2g}",
%              coefficients f"{x:.2f}", percentages f"{100*x:.1f}")
%   \SharpCertBoxSigmaCDLN        S[CD-LN]["sigma_cert1|B3box"]      \SharpCertBoxSigmaCDSU       S[CD-SU]["sigma_cert1|B3box"]
%   \SharpCertSigmaCDLN           S[CD-LN]["sigma_cert1|B3seg"]      \SharpCertSigmaCDSU          S[CD-SU]["sigma_cert1|B3seg"]
%   \SharpCertConeBoxSigmaCDLN    S[CD-LN]["sigma_cert5|B3box"]      \SharpCertConeBoxSigmaCDSU   S[CD-SU]["sigma_cert5|B3box"]
%   \SharpCertConeSigmaCDLN       S[CD-LN]["sigma_cert5|B3seg"]      \SharpCertConeSigmaCDSU      S[CD-SU]["sigma_cert5|B3seg"]
%   \SharpCertBoxScSigmaCDLN      S[CD-LN]["sigma_cert1|B3box_s"]    \SharpCertConeBoxScSigmaCDLN S[CD-LN]["sigma_cert5|B3box_s"]
%   \OldCertBoxScSigmaCDLN        S[CD-LN]["sigma_cert1|B2box_s"]    \OldCertConeBoxScSigmaCDLN   S[CD-LN]["sigma_cert5|B2box_s"]
%   \SharpCertScSigmaCDLN         S[CD-LN]["sigma_cert1|B3seg_s"]    \SharpCertConeScSigmaCDLN    S[CD-LN]["sigma_cert5|B3seg_s"]
%   \SharpEtwoOverBMaxCDLN        S[CD-LN]["maxE2overB|B3box"]       \SharpEtwoOverBMaxCDSU       S[CD-SU]["maxE2overB|B3box"]
%   \SharpEthreeOverBMaxCDLN      S[CD-LN]["maxE3overB3|B3box"]
%   \SharpAold  P["a_old"]   \SharpAnew  P["a_new"]   \SharpAtrue  P["a_true"]   \SharpZmax  P["zmax_median_N2000"]
%   \SharpPredOldKoneLN P["range_aware_LN_old_k1"]  \SharpPredNewKoneLN P["range_aware_LN_new_k1"]
%   \SharpPredOldKfiveLN P["range_aware_LN_old_k5"] \SharpPredNewKfiveLN P["range_aware_LN_new_k5"]
%   \SharpEfourFracPointOne       row sigma=0.1 ["new_B3seg"]     \SharpEfourBoxFracPointOne    row sigma=0.1 ["new_B3box"]
%   \SharpEfourFracPointTwo       row sigma=0.2 ["new_B3seg"]     \SharpEfourBoxFracPointTwo    row sigma=0.2 ["new_B3box"]
%   \SharpEfourBoxScFracPointTwo  row sigma=0.2 ["new_B3box_s"]   \OldEfourBoxScFracPointTwo    row sigma=0.2 ["old_box_s"]
%   \SharpEfourScFracPointTwo     row sigma=0.2 ["new_B3seg_s"]   \SharpEfourScFracPointFour    row sigma=0.4 ["new_B3seg_s"]
%   \OldEfourScFracPointFour      row sigma=0.4 ["old_seg_s"]
%   \SharpCESCertBoxSigmaLN       S[CES-LN]["sigma_cert1|B2box"] (old form, now certified for CES)
%   \SharpCESNewCertBoxSigmaLN    S[CES-LN]["sigma_cert1|B3box"]  \SharpCESNewCertBoxScSigmaLN  S[CES-LN]["sigma_cert1|B3box_s"]
%   \SharpCESTightSmall           max(T["LN_sigma0.01_k3"], T["LN_sigma0.01_k4"])
%   \SharpCESTightLarge           max(T["LN_sigma0.1_k3"],  T["LN_sigma0.1_k4"])
%   \SharpCPUSeconds              int(round(res["meta"]["cpu_seconds"]))
%
% SENTENCES THAT CHANGE ELSEWHERE (suggested wording; macros as above)
%   * Abstract, after "...a fivefold one up to $\sigma=\EoneConeLargestCDLN$)":
%       "; a fourth-order certificate that uses the signed third-moment tensor and the fourth moment, proved
%        here, extends the certified range to $\sigma\le\SharpCertBoxSigmaCDLN$ (fivefold:
%        $\SharpCertConeBoxSigmaCDLN$ instead of $\EoneCertConeBoxSigmaCDLN$), still far from the observed one"
%   * Section "What can be claimed", new row of Table~\ref{tab:claims}:
%       "Fourth-order certificate $(k+1)(|C_3|+B_3)<|Q_2|$ (Prop.~\ref{prop:sharp}); bound holds in every simulated
%        population; certified ($k=1$, LN) up to $\sigma=\SharpCertBoxSigmaCDLN$ from moments and range,
%        $\SharpCertSigmaCDLN$ from micro data; $k=5$: $\SharpCertConeBoxSigmaCDLN$ / $\SharpCertConeSigmaCDLN$ (E1)
%        & proved here; verified"
%     and in the closing paragraph "(valid for $\sigma\le\EoneCertBoxSigmaCDLN$ on log-normal inputs)" ->
%       "(valid for $\sigma\le\EoneCertBoxSigmaCDLN$ on log-normal inputs, $\SharpCertBoxSigmaCDLN$ in its
%        fourth-order form)".
%   * Limitations: "The certified bounds are computed only for Cobb--Douglas, where the derivative suprema have
%       closed forms; for CES a bound would need a numerical supremum, which we did not certify." ->
%       "The derivative suprema have closed forms for Cobb--Douglas; for CES they are enclosed by interval
%        evaluation of the Fa\`a di Bruno formula (theory/check_sharp_certificate.py), which is tight for thin
%        boxes (within a factor $\SharpCESTightSmall$ of a grid maximum at $\sigma=0.01$) and loose for wide
%        ones ($\SharpCESTightLarge$ at $\sigma=0.1$); the CES certificates hold up to
%        $\sigma=\SharpCESNewCertBoxSigmaLN$ (LN, moments and range)."
%     and "because they use absolute third moments where the actual error cancels" -> add "; the fourth-order
%       form removes that cause, and what remains is the growth of the derivative suprema over the input range
%       and the norm inequality (Remark~\ref{rem:sharporders})".
%   * Next steps (a): replace by "(a) The fourth-order certificate (Prop.~\ref{prop:sharp}) exploits the
%       cancellation and changes the order of the bound from $\sigma$ to $\sigma^2$, but its reach is limited by
%       the growth of $\|D^4f\|$ over the coordinate range; a region that controls that growth from moments (or
%       a fifth-order version with the signed fourth-moment tensor) is the next step."
%   * Appendix, after "Moment-and-range form: ...": "Fourth-order certificate: $D^4f$ in the closed form of
%       Lemma~\ref{lem:cd4}, the same corner rule, $B_3$ segmentwise or $\frac N{24}M_4m_4$ on the box, in the
%       Euclidean norm and in the coordinates scaled by $\bar x$; theory/check_sharp_certificate.py reproduces the
%       draws of E1 and E4 and writes theory/sharp_certificate_results.json ($\SharpCPUSeconds$ s of CPU)."
%
% TEST: compiled in a copy of main.tex in the scratchpad with numbers generated from the JSON (pdflatex + bibtex,
% no errors, no undefined references or citations).
% =====================================================================================================

\begin{proposition}[Fourth-order certificate]\label{prop:sharp}
Let $f\in C^4(U)$, $x_1,\dots,x_N\in U$, and let $C_3=\frac N6\langle D^3\!f(\bar x),\mu_3\rangle=\Yh_3-\Yh_2$ be the third-order term, built from the \emph{signed} third central moment tensor. Fix an invertible linear map $S$ of $\R^d$ (below $S=I$ or $S=\operatorname{diag}(\bar x)$), write $\|T\|_S=\|T[S\,\cdot,\dots,S\,\cdot]\|$ for a symmetric $4$-linear form $T$, $M_{4,S}^{(i)}=\sup_{t\in[0,1]}\|D^4\!f(\bar x+th_i)\|_S$, and $m_4^S=\frac1N\sum_i\|S^{-1}h_i\|^4$.
\begin{enumerate}[label=(\alph*),leftmargin=2em]
\item $E_2=C_3+E_3$ with $E_3=Y-\Yh_3=\sum_i\int_0^1\frac{(1-t)^3}{6}\,D^4\!f(\bar x+th_i)[h_i,h_i,h_i,h_i]\,\dd t$.
\item $|E_3|\le\frac1{24}\sum_iM_{4,S}^{(i)}\|S^{-1}h_i\|^4=:B_3\le\frac N{24}M_{4,S}\,m_4^S$ whenever $M_{4,S}\ge\max_iM_{4,S}^{(i)}$, for instance $M_{4,S}=\sup_R\|D^4\!f\|_S$ over a convex region $R\subseteq U$ containing the population. Hence $|E_2|\le|C_3|+B_3$.
\item (Certificate.) $|E_1|\ge|Q_2+C_3|-B_3\ge|Q_2|-|C_3|-B_3$. Consequently, if $(k+1)\big(|C_3|+B_3\big)<|Q_2|$ for some $k\ge1$, or more generally $|Q_2+C_3|-B_3>k\big(|C_3|+B_3\big)$, then $|E_1|>k\big(|C_3|+B_3\big)\ge k|E_2|$.
\item (Data.) With the moment form of (b) the certificate uses $\bar x$, $\Sigma$ (for $Q_2$), $\mu_3$ (the whole tensor, for $C_3$), $m_4^S$ and the region $R$: the \emph{fourth-order moment-and-range certificate}. Unlike $m_3$, $m_4^S=\sum_{j,l}\big(\mu_4^S\big)_{jjll}$ is a contraction of the fourth central moment tensor $\mu_4^S$ of the $S^{-1}h_i$, a polynomial moment. With the segmentwise $B_3$ it needs the individual inputs: the fourth-order \emph{micro-data certificate}.
\item (Comparisons.) For two populations $A$, $B$, $\big|(Y^A-Y^B)-(\Yh_3^A-\Yh_3^B)\big|\le B_3^A+B_3^B$. If $|\Yh_2^A-\Yh_2^B|>\big(|C_3^A|+B_3^A\big)+\big(|C_3^B|+B_3^B\big)$, which is Proposition~\ref{prop:margin} with the bounds of (b), then $|\Yh_3^A-\Yh_3^B|>B_3^A+B_3^B$ and $\Yh_2^A-\Yh_2^B$, $\Yh_3^A-\Yh_3^B$ and $Y^A-Y^B$ have the same sign.
\end{enumerate}
\end{proposition}
\begin{proof}
The segment $t\mapsto\bar x+th_i$, $t\in[0,1]$, lies in the convex set $U$, so $\varphi_i(t)=f(\bar x+th_i)$ is $C^4$ on $[0,1]$ with $\varphi_i^{(r)}(t)=D^r\!f(\bar x+th_i)[h_i,\dots,h_i]$, and Taylor's formula with integral remainder \citep[Ch.~VIII]{dieudonne1960} gives $\varphi_i(1)=\sum_{r=0}^3\varphi_i^{(r)}(0)/r!+\int_0^1\frac{(1-t)^3}{3!}\varphi_i^{(4)}(t)\,\dd t$. Summing over $i$: the terms $r=1$ cancel because $\sum_ih_i=0$; the terms $r=2$ add up to $\frac12\sum_iD^2\!f(\bar x)[h_i,h_i]=Q_2$; and the terms $r=3$ to $\frac16\sum_iD^3\!f(\bar x)[h_i,h_i,h_i]=\frac16\langle D^3\!f(\bar x),\sum_ih_i^{\otimes3}\rangle=C_3$. Hence $Y=Nf(\bar x)+Q_2+C_3+\sum_i\int_0^1\frac{(1-t)^3}{6}\varphi_i^{(4)}(t)\,\dd t$, that is (a), since $\Yh_3=Nf(\bar x)+Q_2+C_3$. For (b), with $v_i=S^{-1}h_i$, $D^4\!f(y)[h_i,\dots,h_i]=D^4\!f(y)[Sv_i,\dots,Sv_i]$ is at most $\|D^4\!f(y)\|_S\|v_i\|^4$ in absolute value, and $\int_0^1\frac{(1-t)^3}{6}\,\dd t=\frac1{24}$; the moment form follows from $M_{4,S}^{(i)}\le M_{4,S}$, and the region $R$ qualifies because, being convex, it contains every segment $[\bar x,x_i]$. For (c), $E_1=E_2+Q_2=Q_2+C_3+E_3$, so $|E_1|\ge|Q_2+C_3|-|E_3|\ge|Q_2+C_3|-B_3$, and $|Q_2+C_3|\ge|Q_2|-|C_3|$; under either hypothesis $|E_1|>k(|C_3|+B_3)$, and $|E_2|\le|C_3|+B_3$ by (b). (d) is a description of what (b) and (c) use, with $\|S^{-1}h\|^4=\sum_{j,l}(S^{-1}h)_j^2(S^{-1}h)_l^2$. For (e), $Y^A-Y^B=(\Yh_3^A-\Yh_3^B)+E_3^A-E_3^B$ by (a); and $|(\Yh_3^A-\Yh_3^B)-(\Yh_2^A-\Yh_2^B)|=|C_3^A-C_3^B|\le|C_3^A|+|C_3^B|$, so the hypothesis gives $|\Yh_3^A-\Yh_3^B|>B_3^A+B_3^B$ with the sign of $\Yh_2^A-\Yh_2^B$, and then the first statement fixes the sign of $Y^A-Y^B$.
\end{proof}

\begin{lemma}[Fourth derivative of a Cobb--Douglas frontier]\label{lem:cd4}
Let $f(x)=A\prod_{m=1}^dx_m^{a_m}$ with $A>0$, $0<a_m<1$, on $\R^d_{>0}$. For a multi-index with multiplicities $e=(e_1,\dots,e_d)$, $\sum_me_m=k$,
\[
\partial^e\!f(x)=A\prod_m(a_m)_{e_m}\,x_m^{a_m-e_m}=f(x)\,\frac{\kappa_e}{\prod_mx_m^{e_m}},\qquad \kappa_e=\prod_m(a_m)_{e_m},
\]
where $(a)_n=a(a-1)\cdots(a-n+1)$ is the falling factorial; for $d=2$, $k=4$ and $(a_1,a_2)=(\alpha,\beta)$ the five distinct coefficients are $\kappa_{(4,0)}=\alpha(\alpha-1)(\alpha-2)(\alpha-3)$, $\kappa_{(3,1)}=\alpha(\alpha-1)(\alpha-2)\beta$, $\kappa_{(2,2)}=\alpha(\alpha-1)\beta(\beta-1)$, $\kappa_{(1,3)}$ and $\kappa_{(0,4)}$ by symmetry. On a box $[\ell,u]\subset\R^d_{>0}$, $|\partial^e\!f|$ is increasing in $x_m$ if $e_m=0$ and decreasing if $e_m\ge1$, so its supremum is attained at the corner $c^e$ with $c^e_m=u_m$ if $e_m=0$ and $c^e_m=\ell_m$ otherwise; and for $S=\operatorname{diag}(s)$, $s\in\R^d_{>0}$,
\[
\sup_{y\in[\ell,u]}\|D^k\!f(y)\|_S\le\Big(\sum_{j_1,\dots,j_k}\big(s^{e}\,|\partial^e\!f(c^e)|\big)^2\Big)^{1/2},\qquad s^e=\prod_ms_m^{e_m},
\]
with $e=e(j_1,\dots,j_k)$ the multiplicities of the index tuple.
\end{lemma}
\begin{proof}
$f$ is a product of one-variable powers, so $\partial^e\!f=A\prod_m\partial_m^{e_m}x_m^{a_m}$ and $\partial^n x^a=(a)_nx^{a-n}$. The exponent $a_m-e_m$ is positive if $e_m=0$ and negative if $e_m\ge1$ (because $a_m<1$), and the coefficient does not depend on $x$, which gives the monotonicity and the corner. Finally $D^k\!f(y)[Sv,\dots,Sv]=\sum_{j_1,\dots,j_k}s^{e}\,\partial^e\!f(y)\,v_{j_1}\cdots v_{j_k}$, which by Cauchy--Schwarz is at most the Frobenius norm of the array $\big(s^e\partial^e\!f(y)\big)$ times $\|v\|^k$, and each entry is bounded by its value at its own corner.
\end{proof}

\begin{remark}[What changes in the order, and what does not]\label{rem:sharporders}
For a centred configuration with vanishing (or $O(\sigma^4)$) third central moments the old bound is $B_2=\Theta(\sigma^3)$ while $|E_2|=O(\sigma^4)$ (Remark~\ref{rem:orders}); the new one is $|C_3|+B_3=O(\sigma^4)$, of the same order as the error it bounds. For log-normal inputs, in the coordinates scaled by $\bar x$ and at leading order with population moments and derivatives at $\bar x$, $B_2/|Q_2|\approx\SharpAold\,\sigma$, $(|C_3|+B_3)/|Q_2|\approx\SharpAnew\,\sigma^2$ and $|E_2|/|Q_2|\approx\SharpAtrue\,\sigma^2$ (constants computed in \path{theory/check_sharp_certificate.py}); the certified dispersion therefore scales like $1/(k+1)$ for the old certificate and like $1/\sqrt{k+1}$ for the new one, which is why the gain is larger for the fivefold certificate. The gain in the exponent is not the whole story at finite $N$: the suprema are taken over the input range, which for log-normal inputs grows like $\bar x\odot e^{\pm\sigma z_{\max}}$ with $z_{\max}\approx\sqrt{2\log N}$, and $\|D^k\!f\|$ grows at the lower corner roughly like $e^{(k-a_m)\sigma z_{\max}}$, faster for $k=4$ than for $k=3$. With the median range of $N=2000$ Gaussian draws ($z_{\max}=\SharpZmax$) the leading-order prediction of the moment-and-range certificate ($k=1$, scaled coordinates) is $\sigma\le\SharpPredNewKoneLN$ against $\SharpPredOldKoneLN$ for the old one, and $\SharpPredNewKfiveLN$ against $\SharpPredOldKfiveLN$ for $k=5$; E1 measures both. The remaining conservatism is a constant factor, from the norm inequality and from that range inflation, not an order; and since the range inflation is stronger for $D^4\!f$, the old bound $B_2$ can be the smaller one at large dispersion, so that $|E_2|\le\min\big(B_2,|C_3|+B_3\big)$ is the bound to use. These orders are heuristic (population moments, leading terms); the certificate itself is exact.
\end{remark}

% ----- [E1-add]: end of the paragraph "Bounds and certificates" of E1 -----
\emph{Fourth-order certificate.} The identity $E_2=C_3+E_3$ of Proposition~\ref{prop:sharp}(a) holds to rounding, and the bound $|E_2|\le|C_3|+B_3$ holds in every replication and cell of both input laws, with $|E_2|/(|C_3|+B_3)$ at most $\SharpEtwoOverBMaxCDLN$ (LN) and $\SharpEtwoOverBMaxCDSU$ (SU) in moment-and-range form ($|E_3|/B_3\le\SharpEthreeOverBMaxCDLN$, LN). In the Euclidean norm used above, the certificate $2(|C_3|+B_3)<|Q_2|$ holds in every replication for $\sigma\le\SharpCertBoxSigmaCDLN$ (LN) and $\sigma\le\SharpCertBoxSigmaCDSU$ (SU) from moments and range, against $\EoneCertBoxSigmaCDLN$ and $\EoneCertBoxSigmaCDSU$ for the third-order one, and for $\sigma\le\SharpCertSigmaCDLN$ and $\SharpCertSigmaCDSU$ in micro-data form (against $\EoneCertSigmaCDLN$ and $\EoneCertSigmaCDSU$); the fivefold improvement is certified for $\sigma\le\SharpCertConeBoxSigmaCDLN$ (LN) and $\SharpCertConeBoxSigmaCDSU$ (SU) from moments and range (against $\EoneCertConeBoxSigmaCDLN$ and $\EoneCertConeBoxSigmaCDSU$) and $\SharpCertConeSigmaCDLN$ and $\SharpCertConeSigmaCDSU$ in micro-data form. Measuring the deviations in the coordinates scaled by $\bar x$ helps both certificates by a comparable factor: from moments and range (LN) the third-order certificate then reaches $\OldCertBoxScSigmaCDLN$ ($k=1$) and $\OldCertConeBoxScSigmaCDLN$ ($k=5$), the fourth-order one $\SharpCertBoxScSigmaCDLN$ and $\SharpCertConeBoxScSigmaCDLN$ ($\SharpCertScSigmaCDLN$ and $\SharpCertConeScSigmaCDLN$ in micro-data form). The signed form of Proposition~\ref{prop:sharp}(c) does not enlarge the certified range in E1 (on SU inputs $C_3=0$). The certified range grows more for $k=5$ than for $k=1$, as Remark~\ref{rem:sharporders} predicts, but remains far below the observed $\sigma=\EoneRatioGtOneSigmaCDLN$ and $\EoneConeLargestCDLN$ (LN).

% ----- [E4-add]: E4, after the sentence that reports \EfourCertBoxFracPointFour -----
With the fourth-order bounds in the margin condition (Proposition~\ref{prop:sharp}(e)) the certified share at $\sigma=0.2$ becomes $\SharpEfourFracPointTwo\%$ (micro data) and $\SharpEfourBoxFracPointTwo\%$ (moments and range), and $\SharpEfourScFracPointTwo\%$ and $\SharpEfourBoxScFracPointTwo\%$ in the coordinates scaled by $\bar x$ (third-order bound, scaled, moments and range: $\OldEfourBoxScFracPointTwo\%$); at $\sigma=0.1$ it is $\SharpEfourFracPointOne\%$ and $\SharpEfourBoxFracPointOne\%$, and at $\sigma=0.4$ at most $\SharpEfourScFracPointFour\%$, less than the third-order bound in scaled coordinates certifies there ($\OldEfourScFracPointFour\%$, micro data): at large dispersion the fourth-order remainder grows faster, and the smaller of the two bounds should be used. No certified comparison is a reversal, as the proposition guarantees.
